\section{UML Class Diagrams}
\label{sec:class}
We show key fields and operations and omit routine getters and constructors.
\texttt{+}, \texttt{-} and \texttt{\#} mean public, private and protected.
Underlining marks static members. Fields marked \texttt{final} cannot be reassigned.
A hollow triangle denotes inheritance, a filled diamond composition, a hollow
diamond aggregation, an arrow association, and a dashed arrow dependency.
Multiplicities describe one running game. State holders are classes of public
static final String constants, not Java enums.

\tikzset{
 rrinherit/.style={-{Triangle[open,length=7pt,width=7pt]}},
 rrcomp/.style={{Diamond[length=7pt,width=6pt]}-{Stealth[length=5pt]}},
 rraggr/.style={{Diamond[open,length=7pt,width=6pt]}-{Stealth[length=5pt]}},
 rrassoc/.style={-{Stealth[length=5pt]}},
 rrdep/.style={dashed,-{Stealth[open,length=5pt]}},
 rrlabel/.style={font=\sffamily\scriptsize,fill=white,inner sep=2pt}}

\subsection{Stage 1 class diagram}

\begin{figure}[p]
\centering
\resizebox{\linewidth}{!}{\begin{tikzpicture}[x=1cm,y=0.85cm,
 rrbox/.style={draw=black!70,fill=classbg,font=\sffamily\small,
 minimum width=3cm,minimum height=0.65cm,align=center},
 rrmul/.style={font=\sffamily\scriptsize,fill=white,inner sep=1.3pt}]
%% boundary & control
\node[rrbox,fill=boundarybg] (cli) at (0,0) {GameCLI};
\node[rrbox,fill=controlbg] (ctrl) at (5,0) {GameController};
\node[rrbox] (main) at (0,2) {Main};
\node[rrbox] (result) at (10,0) {GameResult};
\draw[rrdep] (main.south) -- (cli.north);
\draw[rrassoc] (cli.east) -- node[rrmul,above]{1} (ctrl.west);
\draw[rrdep] (ctrl.east) -- node[rrmul,above]{creates} (result.west);
%% staff hierarchy
\node[rrbox,fill=abstractbg] (staff) at (0,-3) {Staff (abstract)};
\node[rrbox] (waiter) at (-3.5,-5) {Waiter};
\node[rrbox] (chef) at (0,-5) {Chef};
\node[rrbox] (cashier) at (3.5,-5) {Cashier};
\draw[rrinherit] (waiter.north) -- ++(0,0.5) -| (staff.south);
\draw[rrinherit] (chef.north) -- (staff.south);
\draw[rrinherit] (cashier.north) -- ++(0,0.5) -| (staff.south);
%% domain entities
\node[rrbox] (customer) at (8,-3) {Customer};
\node[rrbox] (order) at (12,-3) {Order};
\node[rrbox] (table) at (8,-5) {Table};
\node[rrbox] (menu) at (12,-6) {Menu};
\node[rrbox] (menuitem) at (12,-8) {MenuItem};
\node[rrbox] (payment) at (8,-8) {PaymentRecord};
%% relationships
\draw[rrcomp] (customer.east) -- node[rrmul,above]{0..1} (order.west);
\draw[rrassoc] (table.east) -- ++(0.8,0) |- node[rrmul,right,near start]{0..1} (customer.east);
\draw[rrdep] (customer.south) -- node[rrmul,right]{creates} (payment.north);
\draw[rraggr] (menu.south) -- node[rrmul,right]{3} (menuitem.north);
\draw[rrassoc] (order.south) -- ++(0,-1) -| node[rrmul,left,near end]{1} (menuitem.north east);
%% ownership from controller
\draw[rrcomp] (ctrl.south) -- ++(0,-0.5) -| node[rrmul,left,near end]{1} (waiter.north west);
\draw[rrcomp] (ctrl.south) -- ++(0,-0.5) -| (chef.north);
\draw[rrcomp] (ctrl.south) -- ++(0,-0.5) -| node[rrmul,right,near end]{1} (cashier.north east);
\draw[rrcomp] (ctrl.south) -- ++(0,-0.5) -| node[rrmul,left,near end]{0..6} (customer.north);
\draw[rrcomp] (ctrl.south) -- ++(0,-0.5) -| node[rrmul,left,near end]{2} (table.north);
\draw[rrcomp] (ctrl.south) -- ++(0,-0.5) -| node[rrmul,right,near end]{1} (menu.north);
%% constants
\node[rrbox,fill=enumbg,font=\sffamily\footnotesize] (cs) at (4,-8) {CustomerStatus};
\node[rrbox,fill=enumbg,font=\sffamily\footnotesize] (os) at (4,-9.5) {OrderStatus};
\node[rrbox,fill=enumbg,font=\sffamily\footnotesize] (tp) at (4,-11) {TurnPhase};
\draw[rrdep] (customer.south west) -- (cs.east);
\draw[rrdep] (order.south) -- ++(0,-3) -| (os.east);
\draw[rrdep] (ctrl.south) -- ++(0,-8) -| (tp.north);
\end{tikzpicture}}
\caption{Stage~1 class diagram. Three concrete staff roles inherit the one-task availability
from abstract Staff. Customer owns at most one Order; GameController owns all domain objects.
The three String-constant holders act as state enumerations.
17 classes total.}
\label{fig:class-stage1}
\end{figure}
\clearpage

\subsection{Stage 2 class diagram}
Stage~2 adds polymorphic customer types and a schedule that constructs
them. Three new classes (shaded darker) extend the Stage~1 design without
changing existing interfaces.

\begin{figure}[p]
\centering
\resizebox{\linewidth}{!}{\begin{tikzpicture}[x=1cm,y=0.85cm,
 rrbox/.style={draw=black!70,fill=classbg,font=\sffamily\small,
 minimum width=3cm,minimum height=0.65cm,align=center},
 rrnew/.style={draw=black!70,fill=orange!18,font=\sffamily\small,
 minimum width=3cm,minimum height=0.65cm,align=center},
 rrmul/.style={font=\sffamily\scriptsize,fill=white,inner sep=1.3pt}]
%% boundary & control
\node[rrbox,fill=boundarybg] (cli) at (0,0) {GameCLI};
\node[rrbox,fill=controlbg] (ctrl) at (5,0) {GameController};
\node[rrbox] (main) at (0,2) {Main};
\node[rrbox] (result) at (10,0) {GameResult};
\draw[rrdep] (main.south) -- (cli.north);
\draw[rrassoc] (cli.east) -- node[rrmul,above]{1} (ctrl.west);
\draw[rrdep] (ctrl.east) -- node[rrmul,above]{creates} (result.west);
%% NEW: ArrivalSchedule
\node[rrnew] (arr) at (10,2) {ArrivalSchedule \textit{(new)}};
\draw[rrcomp] (ctrl.north east) -- ++(0,0.5) -| node[rrmul,right,near end]{1} (arr.south);
%% staff hierarchy
\node[rrbox,fill=abstractbg] (staff) at (0,-3) {Staff (abstract)};
\node[rrbox] (waiter) at (-3.5,-5) {Waiter};
\node[rrbox] (chef) at (0,-5) {Chef};
\node[rrbox] (cashier) at (3.5,-5) {Cashier};
\draw[rrinherit] (waiter.north) -- ++(0,0.5) -| (staff.south);
\draw[rrinherit] (chef.north) -- (staff.south);
\draw[rrinherit] (cashier.north) -- ++(0,0.5) -| (staff.south);
%% domain entities
\node[rrbox] (customer) at (8,-3) {Customer};
\node[rrbox] (order) at (12,-3) {Order};
\node[rrbox] (table) at (8,-5) {Table};
\node[rrbox] (menu) at (12,-6) {Menu};
\node[rrbox] (menuitem) at (12,-8) {MenuItem};
\node[rrbox] (payment) at (8,-8) {PaymentRecord};
%% NEW: VIP and Critic
\node[rrnew] (vip) at (6.5,-6.5) {VIPCustomer \textit{(new)}};
\node[rrnew] (critic) at (10,-6.5) {CriticCustomer \textit{(new)}};
\draw[rrinherit] (vip.north) -- ++(0,0.5) -| (customer.south);
\draw[rrinherit] (critic.north) -- ++(0,0.5) -| (customer.south east);
%% relationships
\draw[rrcomp] (customer.east) -- node[rrmul,above]{0..1} (order.west);
\draw[rrassoc] (table.east) -- ++(0.8,0) |- node[rrmul,right,near start]{0..1} (customer.east);
\draw[rrdep] (customer.south west) -- ++(0,-3.5) -| (payment.south);
\draw[rraggr] (menu.south) -- node[rrmul,right]{3} (menuitem.north);
\draw[rrassoc] (order.south) -- ++(0,-1) -| node[rrmul,left,near end]{1} (menuitem.north east);
%% ownership from controller
\draw[rrcomp] (ctrl.south) -- ++(0,-0.5) -| node[rrmul,left,near end]{1} (waiter.north west);
\draw[rrcomp] (ctrl.south) -- ++(0,-0.5) -| (chef.north);
\draw[rrcomp] (ctrl.south) -- ++(0,-0.5) -| node[rrmul,right,near end]{1} (cashier.north east);
\draw[rrcomp] (ctrl.south) -- ++(0,-0.5) -| node[rrmul,left,near end]{0..6} (customer.north);
\draw[rrcomp] (ctrl.south) -- ++(0,-0.5) -| node[rrmul,left,near end]{2} (table.north);
\draw[rrcomp] (ctrl.south) -- ++(0,-0.5) -| node[rrmul,right,near end]{1} (menu.north);
%% constants
\node[rrbox,fill=enumbg,font=\sffamily\footnotesize] (cs) at (4,-9) {CustomerStatus};
\node[rrbox,fill=enumbg,font=\sffamily\footnotesize] (os) at (4,-10.5) {OrderStatus};
\node[rrbox,fill=enumbg,font=\sffamily\footnotesize] (tp) at (4,-12) {TurnPhase};
\draw[rrdep] (customer.south west) -- ++(0,-4) -| (cs.east);
\draw[rrdep] (order.south) -- ++(0,-3.5) -| (os.east);
\draw[rrdep] (ctrl.south) -- ++(0,-9) -| (tp.north);
\end{tikzpicture}}
\caption{Stage~2 class diagram. Shaded boxes are new. ArrivalSchedule constructs VIP and
Critic subclasses via a factory method. VIPCustomer overrides patience loss and adds a
10\% bill discount; CriticCustomer overrides meal selection to prefer the most expensive
well-prepared dish. Existing Customer references dispatch the differing rules via polymorphism.
20 classes total (+3 from Stage~1).}
\label{fig:class-stage2}
\end{figure}
\clearpage

\subsection{Stage 3 class diagrams (detailed views)}
The following complementary views describe the final Stage~3 implementation.
They share class names across views; repeated boxes refer to the same classes.
Figure~\ref{fig:class-relationships} gathers the ownership, associations and state dependencies.

\tikzset{
 rrclass/.style={rectangle split,rectangle split parts=3,
 rectangle split part fill={classbg,white,white},draw=black!70,
 font=\ttfamily\scriptsize,text width=6.2cm,align=left,inner sep=4pt,
 rectangle split part align={center,left,left}}}

\begin{figure}[p]
\centering
\resizebox{\linewidth}{!}{\begin{tikzpicture}
\node[rrclass,anchor=north west] (main) at (0,0) {\textbf{Main}
\nodepart{second}(no fields shown)
\nodepart{third}\underline{+ main(args: String[]): void}};
\node[rrclass,anchor=north west] (cli) at (0,-2.5) {\textbf{GameCLI}
\nodepart{second}- game: GameController {final}\\
- input: Scanner {final}\\
- inputEnded: boolean
\nodepart{third}+ run(): void\\
- beginTurn(): void\\
- choose(line: String, role: String): void\\
- displayState(): void};
\node[rrclass,anchor=north west] (ctrl) at (7,0) {\textbf{GameController}
\nodepart{second}\underline{+ STAGE: int = 3}\\
\underline{+ MAX\_TURNS: int = 18}\\
- menu: Menu {final}\\
- customers: ArrayList$<$Customer$>$ {final}\\
- tables: ArrayList$<$Table$>$ {final}\\
- payments: ArrayList$<$PaymentRecord$>$ {final}\\
- arrivals: ArrivalSchedule {final}\\
- happyHour: HappyHour {final}\\
- host: Host {final}\\
- waiter: Waiter {final}\\
- chef: Chef {final}\\
- cashier: Cashier {final}\\
- turn: int\\
- phase: String
\nodepart{third}+ beginTurn(): void\\
+ beginTurn(activateHappyHour: boolean): void\\
+ seatWithHost(id: int): void\\
+ waitHost(): void\\
+ seatCustomer(id: int): void\\
+ takeOrder(id: int): void\\
+ serveCustomer(id: int): void\\
+ waitWaiter(): void\\
+ prepareOrder(id: int): void\\
+ waitChef(): void\\
+ finishTurn(): void\\
+ getResult(): GameResult\\
- requirePhase(expected: String): void};
\node[rrclass,anchor=north west] (arr) at (0,-6.6) {\textbf{ArrivalSchedule}
\nodepart{second}(no fields shown)
\nodepart{third}+ arrivalAt(turn: int): Customer};
\node[rrclass,anchor=north west] (result) at (0,-8.7) {\textbf{GameResult}
\nodepart{second}- paidCustomers: int {final}\\
- revenueCents: int {final}\\
- satisfactionTotal: int {final}
\nodepart{third}+ isVictory(): boolean};
\draw[rrdep] (main.south) -- (cli.north);
\draw[rrassoc] (cli.east) -- node[rrlabel,above]{1} (ctrl.west);
\draw[rrcomp] (ctrl.west |- arr.east) -- node[rrlabel,above]{1} (arr.east);
\draw[rrdep] (ctrl.west |- result.east) -- node[rrlabel,above]{creates} (result.east);
\end{tikzpicture}}
\caption{Boundary and turn control. Main constructs the CLI and controller; GameCLI requests actions and catches invalid choices. GameResult checks all three thresholds using unrounded payment totals.}
\label{fig:class-full}
\end{figure}
\clearpage
\begin{figure}[p]
\centering
\resizebox{\linewidth}{!}{\begin{tikzpicture}
\node[rrclass,anchor=north west] (customer) at (0,0) {\textbf{Customer}
\nodepart{second}- id: int {final}\\
- arrivalTurn: int {final}\\
- satisfaction: int\\
- status: String\\
- order: Order\\
- servedTurn: int\\
- readyToPayTurn: int
\nodepart{third}+ chooseMeal(menu: Menu): MenuItem\\
+ seat(): void\\
+ placeOrder(newOrder: Order): void\\
+ serve(turn: int): void\\
+ finishEating(turn: int): void\\
+ canPay(turn: int): boolean\\
+ pay(turn: int): PaymentRecord\\
+ losePatience(): void\\
+ recoverSatisfaction(amount: int): void\\
+ getWaitingLoss(): int\\
+ discountedBill(baseCents: int): int};
\node[rrclass,anchor=north west] (order) at (7,0) {\textbf{Order}
\nodepart{second}- item: MenuItem {final}\\
- billCents: int {final}\\
- remainingUnits: int\\
- readyTurn: int\\
- status: String
\nodepart{third}+ prepare(turn: int): void\\
+ serve(turn: int): void\\
+ markPaid(): void\\
+ cancel(): void};
\node[rrclass,anchor=north west] (pay) at (7,-5.9) {\textbf{PaymentRecord}
\nodepart{second}- customerId: int {final}\\
- billCents: int {final}\\
- satisfaction: int {final}\\
- turn: int {final}
\nodepart{third}+ getCustomerId(): int\\
+ getBillCents(): int\\
+ getSatisfaction(): int\\
+ getTurn(): int};
\node[rrclass,anchor=north west] (vip) at (0,-10) {\textbf{VIPCustomer}
\nodepart{second}(no fields shown)
\nodepart{third}+ getWaitingLoss(): int\\
+ getType(): String\\
+ discountedBill(baseCents: int): int};
\node[rrclass,anchor=north west] (critic) at (7,-10) {\textbf{CriticCustomer}
\nodepart{second}(no fields shown)
\nodepart{third}+ getWaitingLoss(): int\\
+ getType(): String\\
+ chooseMeal(menu: Menu): MenuItem};
\draw[rrcomp] (customer.east) -- node[rrlabel,above]{1 owns 0..1} (order.west);
\draw[rrdep] (customer.east |- pay.west) -- node[rrlabel,above]{creates at payment} (pay.west);
\draw[rrinherit] (vip.north) -- (customer.south);
\draw[rrinherit] (critic.north) -- ++(0,0.5) -| (customer.south east);
\end{tikzpicture}}
\caption{Customer and order lifecycle. Encapsulation: transition methods guard private state and timestamps. Composition: each Customer owns at most one Order. Polymorphism: Customer is the concrete Regular type; VIP and Critic override selected behaviour through Customer references. PaymentRecord preserves the bill and satisfaction at payment.}
\label{fig:class-customers}
\end{figure}
\clearpage

\begin{figure}[p]
\centering
\resizebox{\linewidth}{!}{\begin{tikzpicture}
\node[rrclass,anchor=north west] (staff) at (0,0) {\textbf{Staff (abstract)}
\nodepart{second}- available: boolean
\nodepart{third}+ startTurn(): void {final}\\
+ isAvailable(): boolean {final}\\
\# requireAvailable(): void {final}\\
\# completeTask(): void {final}};
\node[rrclass,anchor=north west] (host) at (0,-4) {\textbf{Host}
\nodepart{second}(no fields shown)
\nodepart{third}+ seat(customer: Customer,\\
  table: Table): void};
\node[rrclass,anchor=north west] (waiter) at (7,-4) {\textbf{Waiter}
\nodepart{second}- pricing: PricingService {final}
\nodepart{third}+ seat(customer: Customer,\\
  table: Table): void\\
+ takeOrder(customer: Customer,\\
  menu: Menu): void\\
+ serve(customer: Customer, turn: int): void};
\node[rrclass,anchor=north west] (chef) at (0,-7.5) {\textbf{Chef}
\nodepart{second}(no fields shown)
\nodepart{third}+ prepare(customer: Customer,\\
  turn: int): void};
\node[rrclass,anchor=north west] (cash) at (0,-10.4) {\textbf{Cashier}
\nodepart{second}(no fields shown)
\nodepart{third}+ collectOldest(customers:\\
  ArrayList$<$Customer$>$, turn: int):\\
  PaymentRecord};
\node[rrclass,anchor=north west] (price) at (7,-9) {\textbf{PricingService}
\nodepart{second}- happyHour: HappyHour {final}
\nodepart{third}+ quote(customer: Customer,\\
  item: MenuItem): int};
\node[rrclass,anchor=north west] (hh) at (7,-12.6) {\textbf{HappyHour}
\nodepart{second}- currentTurn: int\\
- activationTurn: int
\nodepart{third}+ startTurn(turn: int, activate: boolean): void\\
+ hasBeenUsed(): boolean\\
+ isActive(): boolean\\
+ getRemainingTurns(): int};
\draw[rrinherit] (host.north) -- (staff.south);
\draw[rrinherit] (waiter.north) -- ++(0,0.5) |- (staff.east);
\draw[rrinherit] (chef.west) -- ++(-0.5,0) |- (staff.west);
\draw[rrinherit] (cash.west) -- ++(-0.8,0) |- (staff.west);
\draw[rrassoc] (waiter.south) -- node[rrlabel,right]{1} (price.north);
\draw[rrassoc] (price.south) -- node[rrlabel,right]{1} (hh.north);
\end{tikzpicture}}
\caption{Staff and pricing. Staff shares one-task availability without declaring an unrelated generic work operation. Host and Waiter depend on Table; Chef depends on Order; Cashier depends on Customer and creates no payment until a customer is eligible. PricingService references the same HappyHour owned by the controller.}
\label{fig:class-staff}
\end{figure}
\clearpage

\begin{figure}[p]
\centering
\resizebox{\linewidth}{!}{\begin{tikzpicture}
\node[rrclass,anchor=north west] (menu) at (0,0) {\textbf{Menu}
\nodepart{second}- items: ArrayList$<$MenuItem$>$ {final}
\nodepart{third}\underline{+ withCombo(): Menu}\\
+ getItems(): ArrayList$<$MenuItem$>$};
\node[rrclass,anchor=north west] (item) at (7,0) {\textbf{MenuItem}
\nodepart{second}\underline{+ SALAD: MenuItem {final}}\\
\underline{+ BURGER: MenuItem {final}}\\
\underline{+ PASTA: MenuItem {final}}\\
- name: String {final}\\
- priceCents: int {final}\\
- preparationUnits: int {final}
\nodepart{third}+ getName(): String\\
+ getPriceCents(): int\\
+ getPreparationUnits(): int};
\node[rrclass,anchor=north west] (combo) at (7,-5.6) {\textbf{ComboMeal}
\nodepart{second}- first: MenuItem {final}\\
- second: MenuItem {final}
\nodepart{third}+ getComponents(): ArrayList$<$MenuItem$>$};
\node[rrclass,anchor=north west] (table) at (0,-5.6) {\textbf{Table}
\nodepart{second}- id: int {final}\\
- occupant: Customer
\nodepart{third}+ seat(customer: Customer): void\\
+ releaseIfFinished(): void\\
+ isFree(): boolean};
\node[rrclass,anchor=north west] (cs) at (0,-10) {\textbf{CustomerStatus}
\nodepart{second}\underline{+ WAITING, SEATED, EATING: String}\\
\underline{+ READY\_TO\_PAY, PAID, LEFT: String}
\nodepart{third}(no operations shown)};
\node[rrclass,anchor=north west] (os) at (7,-10) {\textbf{OrderStatus}
\nodepart{second}\underline{+ TAKEN, PREPARING, READY: String}\\
\underline{+ SERVED, PAID, CANCELLED: String}
\nodepart{third}(no operations shown)};
\node[rrclass,anchor=north west] (phase) at (0,-13) {\textbf{TurnPhase}
\nodepart{second}\underline{+ BETWEEN\_TURNS, HOST, WAITER: String}\\
\underline{+ CHEF, END\_OF\_TURN, FINISHED: String}
\nodepart{third}(no operations shown)};
\draw[rraggr] (menu.east) -- node[rrlabel,above]{1 holds 4} (item.west);
\draw[rrinherit] (combo.north) -- (item.south);
\draw[rraggr] (combo.east) -- ++(0.7,0) |- node[rrlabel,right]{2 components} (item.east);
\end{tikzpicture}}
\caption{Menu, table and String constant holders. Menu items and combo components are shared immutable references, so hollow diamonds indicate aggregation. Each Table references 0..1 Customer; each Customer occupies 0..1 Table. An Order references exactly one MenuItem. Customer, Order and GameController depend on their respective status/phase constant classes. All constants shown are public static final; grouped names denote separate fields.}
\label{fig:class-menu}
\end{figure}
\clearpage

\begin{landscape}
\begin{figure}[p]
\centering
\resizebox{0.98\linewidth}{!}{%
\begin{tikzpicture}[x=1cm,y=0.85cm,
 rrbox/.style={draw=black!70,fill=classbg,font=\sffamily\small,
 minimum width=3cm,minimum height=0.65cm,align=center},
 rrmul/.style={font=\sffamily\scriptsize,fill=white,inner sep=1.3pt}]
\node[font=\sffamily\bfseries] at (3,1) {Game-session ownership};
\node[rrbox,fill=controlbg] (owner) at (0,-5.5) {GameController};
\draw[{Diamond[length=7pt,width=6pt]}-] (owner.east) -- (2.5,-5.5);
\node[rrmul,above] at (1.9,-5.5) {1};
\draw (2.5,0) -- (2.5,-11);
\node[rrbox] (owned0) at (5.1,0.0) {Table};
\draw[rrassoc] (2.5,0.0) -- node[rrmul,above,near end] {2} (owned0.west);
\node[rrbox] (owned1) at (5.1,-1.2) {Customer};
\draw[rrassoc] (2.5,-1.2) -- node[rrmul,above,near end] {0..6} (owned1.west);
\node[rrbox] (owned2) at (5.1,-2.4) {PaymentRecord};
\draw[rrassoc] (2.5,-2.4) -- node[rrmul,above,near end] {0..5} (owned2.west);
\node[rrbox] (owned3) at (5.1,-3.6) {Menu};
\draw[rrassoc] (2.5,-3.6) -- node[rrmul,above,near end] {1} (owned3.west);
\node[rrbox] (owned4) at (5.1,-4.8) {ArrivalSchedule};
\draw[rrassoc] (2.5,-4.8) -- node[rrmul,above,near end] {1} (owned4.west);
\node[rrbox] (owned5) at (5.1,-6.0) {HappyHour};
\draw[rrassoc] (2.5,-6.0) -- node[rrmul,above,near end] {1} (owned5.west);
\node[rrbox] (owned6) at (5.1,-7.2) {Host};
\draw[rrassoc] (2.5,-7.2) -- node[rrmul,above,near end] {1} (owned6.west);
\node[rrbox] (owned7) at (5.1,-8.4) {Waiter};
\draw[rrassoc] (2.5,-8.4) -- node[rrmul,above,near end] {1} (owned7.west);
\node[rrbox] (owned8) at (5.1,-9.6) {Chef};
\draw[rrassoc] (2.5,-9.6) -- node[rrmul,above,near end] {1} (owned8.west);
\node[rrbox] (owned9) at (5.1,-10.8) {Cashier};
\draw[rrassoc] (2.5,-10.8) -- node[rrmul,above,near end] {1} (owned9.west);

\node[font=\sffamily\bfseries] at (15,1) {Domain and state relationships};
\node[rrbox] (rt) at (10,0) {Table};
\node[rrbox] (rc) at (15,0) {Customer};
\node[rrbox,fill=enumbg] (rcs) at (20,0) {CustomerStatus};
\node[rrbox] (ro) at (15,-2.5) {Order};
\node[rrbox,fill=enumbg] (ros) at (20,-2.5) {OrderStatus};
\node[rrbox] (rm) at (10,-5) {Menu};
\node[rrbox] (ri) at (15,-5) {MenuItem};
\node[rrbox] (rcombo) at (20,-5) {ComboMeal};
\node[rrbox] (rp) at (10,-8) {PricingService};
\node[rrbox] (rh) at (15,-8) {HappyHour};
\node[rrbox] (rw) at (10,-10.5) {Waiter};
\node[rrbox,fill=controlbg] (rg) at (15,-10.5) {GameController};
\node[rrbox,fill=enumbg] (rphase) at (20,-10.5) {TurnPhase};
\draw[rrassoc] (rt.east) -- node[rrmul,above]{occupant}
 node[rrmul,below,near start]{0..1} node[rrmul,below,near end]{0..1} (rc.west);
\draw[rrcomp] (rc.south) -- node[rrmul,right,near start]{1}
 node[rrmul,right,near end]{0..1} (ro.north);
\draw[rrassoc] (ro.south) -- node[rrmul,right,near start]{0..*}
 node[rrmul,right,near end]{1} (ri.north);
\draw[rraggr] (rm.east) -- node[rrmul,above]{items}
 node[rrmul,below,near start]{0..*} node[rrmul,below,near end]{3..4} (ri.west);
\draw[rraggr] (rcombo.west) -- node[rrmul,above]{components}
 node[rrmul,below,near start]{0..*} node[rrmul,below,near end]{2} (ri.east);
\draw[rrassoc] (rp.east) -- node[rrmul,above]{happyHour}
 node[rrmul,below,near start]{1} node[rrmul,below,near end]{1} (rh.west);
\draw[rrassoc] (rw.north) -- node[rrmul,left,near start]{1}
 node[rrmul,left,near end]{1} (rp.south);
\draw[rrdep] (rc.east) -- (rcs.west);
\draw[rrdep] (ro.east) -- (ros.west);
\draw[rrdep] (rg.east) -- (rphase.west);
\end{tikzpicture}}
\caption{Relationship overview. The shared composition trunk on the left means that each
listed object is owned by one GameController; numbers beside the destinations are
per-game counts. On the right, endpoint numbers are association multiplicities.
Menu and ComboMeal share immutable menu entries (hollow diamonds); Customer owns
its Order (filled diamond). State dependencies point to String constant classes,
not enums. Repeated names denote the same classes as in the detailed views.}
\label{fig:class-relationships}
\end{figure}
\end{landscape}
\clearpage
\noindent\textbf{Implementation notes.} Menu returns a copy of its list, but the immutable
menu entries are shared. GameController similarly returns container copies for
customers, tables and payments. Customer, Order and Table enforce their own
valid transitions; the CLI receives exceptions and retries the same role.
The model deliberately uses ordinary classes, ArrayLists and loops rather than
additional interfaces or a framework.
